\section{Phân tích yêu cầu}
    Mục tiêu cốt lõi là thiết kế một IoT Gateway giải quyết bài toán linh hoạt và khả cấu hình nhằm tối ưu chi phí và thời gian triển khai hệ thống cho đa dạng các nhu cầu sử dụng.

    \subsection{Yêu cầu chức năng}
        Hệ thống Gateway IoT cần đáp ứng các nhóm chức năng chính sau đây nhằm đảm bảo tính linh hoạt và khả năng xử lý dữ liệu tin cậy:

        \begin{itemize}
            \item \textbf{Tính mô-đun hóa phần cứng:}
            \begin{itemize}
                \item Thiết kế theo kiến trúc mô-đun, hỗ trợ kết nối linh hoạt và tương thích với đa dạng các loại mô-đun.
                \item Cung cấp cơ chế thay thế, nâng cấp mô-đun tùy theo mục đích ứng dụng mà không yêu cầu tái thiết kế bo mạch chủ (Baseboard).
                \item Hỗ trợ các chuẩn giao tiếp ngoại vi phổ thông: UART, SPI, I2C, USB, \dots
            \end{itemize}
            
            \item \textbf{Khả năng thu thập và xử lý dữ liệu:}
            \begin{itemize}
                \item Hỗ trợ đa giao tiếp và đa nền tảng: Gateway phải có khả năng tiếp nhận dữ liệu đồng thời từ nhiều nhóm thiết bị đầu cuối hoạt động trên các chuẩn giao tiếp khác nhau như Zigbee, LoRa, Bluetooth, và các giao thức công nghiệp (RS-485/\ldots).
                \item Xử lý song song và biệt lập: Hệ thống phải triển khai các tiến trình xử lý độc lập cho từng giao thức truyền thông, đảm bảo việc thu thập dữ liệu từ một luồng không gây tắc nghẽn hoặc ảnh hưởng đến các luồng dữ liệu khác.
                \item Tiền xử lý và Chuẩn hóa dữ liệu tại biên: Thực hiện kiểm chứng định dạng, kiểm tra tính toàn vẹn (Checksum/CRC) và loại bỏ các khung tin lỗi ngay tại tầng LAN.
                \item Cơ chế đệm và đảm bảo an toàn dữ liệu: Sử dụng hàng đợi (Queue) và bộ nhớ đệm để lưu trữ dữ liệu tạm thời, đảm bảo không mất mát thông tin trong trường hợp kết nối mạng diện rộng (WAN) bị gián đoạn.
                \item Chuyển giao dữ liệu linh hoạt: Hỗ trợ đóng gói và chuyển phát dữ liệu lên các máy chủ quản lý thông qua các giao thức ứng dụng phổ biến như MQTT, HTTP hoặc CoAP tùy theo yêu cầu hệ thống.
            \end{itemize}
            
            \item \textbf{Khả năng khả cấu hình:}
            \begin{itemize}
                \item Cấu hình không xâm lấn (Non-intrusive Configuration): Cho phép tinh chỉnh toàn bộ tham số vận hành, mạng và giao thức mà không cần can thiệp vào mã nguồn hoặc biên dịch lại Firmware.
                \item Định nghĩa Module bằng phần mềm (Software-Defined Module): Hỗ trợ cơ chế cấu hình động dựa trên tệp mô tả (JSON) để đặc tả chi tiết thông số vật lý (Baudrate, Parity), tập lệnh điều khiển (Command Template), trình tự điều khiển chân logic (GPIO Sequence) và các tham số thời gian (Timeout) của module.
                \item Đa dạng phương thức tương tác: Cung cấp giao diện cấu hình trực quan qua ứng dụng máy tính (PC Tools) và giao diện Web (Web Config Portal) nhúng trực tiếp trên thiết bị.
                \item Quản lý cấu hình thông minh: Hệ thống phải tự động lưu trữ cấu hình vào bộ nhớ không khả biến (NVS) và có khả năng tự phục hồi trạng thái vận hành sau khi khởi động lại.
                \item Linh hoạt trong triển khai thực địa: Hỗ trợ chế độ điểm truy cập (AP Mode) với Captive Portal để cấu hình mạng ban đầu và chế độ trạm (STA Mode) để vận hành trong mạng nội bộ.
            \end{itemize}
        \end{itemize}

    \subsection{Yêu cầu phi chức năng}
        IoT Gateway phải thỏa mãn các tiêu chuẩn vận hành để đáp ứng điều kiện triển khai thực tế:
        \begin{itemize}
            \item \textbf{Độ tin cậy và tính ổn định:}
            \begin{itemize}
                \item Đảm bảo khả năng vận hành liên tục 24/7 trong thời gian dài.
                \item Tích hợp cơ chế giám sát phần cứng và phần mềm (Watchdog Timer) để tự động phục hồi khi xảy ra hiện tượng treo hệ thống hoặc xung đột tiến trình.
                \item Chủ động phát hiện và xử lý có kiểm soát các lỗi ngoại vi (mất Internet, lỗi cảm biến, lỗi mô-đun giao tiếp) thông qua hệ thống bản ghi (Log) và chuyển đổi về trạng thái an toàn.
                \item Duy trì tính toàn vẹn của dữ liệu bằng bộ đệm nội bộ, hỗ trợ lưu trữ và gửi bù dữ liệu ngay khi kết nối mạng được khôi phục (trong khoảng 200 đến 300 ms kể từ khi có kết nối lại với server).
                \item Hỗ trợ cập nhật phần mềm từ xa (FOTA) an toàn cho nhân xử lý với cơ chế kiểm tra tính toàn vẹn, đảm bảo không ghi đè lên Firmware hiện tại khi chưa xác thực thành công bản cập nhật.
                \item Thiết kế phần cứng có khả năng chống chịu tác động môi trường: bảo vệ chống tĩnh điện (ESD), lọc nhiễu điện từ (EMI).
            \end{itemize}
            
            \item \textbf{Nguồn điện và năng lượng:}
            \begin{itemize}
                \item Hệ thống nguồn yêu cầu độ ổn định cao, tích hợp các mạch bảo vệ quá áp (OVP), bảo vệ ngắn mạch (SCP) và chống ngược cực.
                \item Tối ưu hóa mức tiêu thụ năng lượng thông qua các chế độ tiết kiệm điện (Sleep mode) của vi điều khiển, phù hợp cho các ứng dụng vận hành bằng pin hoặc năng lượng tái tạo.
            \end{itemize}
            
            \item \textbf{Hiệu năng:}
            \begin{itemize}
                \item Năng lực xử lý trung tâm đảm bảo độ trễ thấp (dưới 1 ms) trong việc thu thập và chuyển giao dữ liệu từ tất cả các mô-đun kết nối đồng thời.
                \item Tốc độ phản hồi giao diện cấu hình và thời gian khởi động hệ thống tối ưu (khoảng 15 giây cho WiFi và 40 giây cho LTE), đảm bảo không gây tắc nghẽn hàng đợi dữ liệu khi mật độ thiết bị đầu cuối tăng cao (có thể kết nối từ 30--40 thiết bị).
                \item Thiết kế có thể xử lý tối đa kết nối với khoảng 30--40 thiết bị, tuy nhiên còn phụ thuộc vào module sử dụng cho giao thức.
            \end{itemize}
            
            \item \textbf{Tính khả dụng và dễ sử dụng:}
            \begin{itemize}
                \item Thiết kế cơ khí hỗ trợ tháo lắp mô-đun nhanh chóng, giảm thiểu yêu cầu về công cụ chuyên dụng.
                \item Giao diện phần mềm thân thiện, lược bỏ các thao tác lập trình phức tạp, cho phép người dùng cấu hình thông qua các tùy chọn có sẵn.
                \item Hệ thống chỉ thị trực quan qua đèn LED (trạng thái nguồn, kết nối LAN/WAN, trạng thái Server, lỗi hệ thống). Duy trì nhật ký sự kiện chi tiết và cung cấp bộ tài liệu hướng dẫn vận hành, xử lý sự cố (Troubleshooting) toàn diện.
            \end{itemize}
        \end{itemize}

    \subsection{Các ràng buộc thiết kế}
        \begin{itemize}
            \item \textbf{Ràng buộc về chi phí:} Cần đảm bảo chi phí sản xuất và triển khai gateway trên thị trường có khả năng cạnh tranh, tối ưu hóa lựa chọn phần cứng và phần mềm sao cho tiết kiệm mà vẫn đáp ứng được yêu cầu hệ thống.
            \item \textbf{Ràng buộc về công nghệ:} Cần phải sử dụng các công nghệ phổ biến và đáng tin cậy, đảm bảo khả năng tương thích với các thiết bị hiện có, cũng như dễ dàng tích hợp vào các nền tảng cloud hoặc hệ thống IoT.
            \item \textbf{Ràng buộc về hiệu suất:} Gateway phải đảm bảo tốc độ truyền dữ liệu, khả năng xử lý song song và đáp ứng các yêu cầu về độ trễ thấp trong quá trình thu thập và truyền tải dữ liệu.
            \item \textbf{Ràng buộc về khả năng mở rộng:} Hệ thống phải được thiết kế sao cho dễ dàng mở rộng hoặc thay thế các module mà không làm gián đoạn hoạt động của hệ thống, hỗ trợ thêm giao thức hoặc phần cứng khi cần thiết.
        \end{itemize}